\documentclass[tikz,border=10pt]{standalone}
\usepackage[no-math]{fontspec}
\usepackage[brazil]{babel}
\usepackage{amsmath,amssymb}
\IfFileExists{C:/Windows/Fonts/Lato_Regular.ttf}{%
  \setmainfont{Lato}[
    Path = C:/Windows/Fonts/,
    UprightFont = Lato_Regular.ttf,
    BoldFont = Lato_Bold.ttf,
    ItalicFont = Lato_Italic.ttf,
    BoldItalicFont = Lato_BoldItalic.ttf
  ]
  \setsansfont{Lato}[
    Path = C:/Windows/Fonts/,
    UprightFont = Lato_Regular.ttf,
    BoldFont = Lato_Bold.ttf,
    ItalicFont = Lato_Italic.ttf,
    BoldItalicFont = Lato_BoldItalic.ttf
  ]
}{%
  \RequirePackage[default]{lato}
}
\renewcommand{\familydefault}{\sfdefault}

\usepackage{xcolor}
\usepackage{tikz}
\usepackage{pgfplots}
\pgfplotsset{compat=1.18}
\usetikzlibrary{arrows.meta,positioning,shapes.geometric,fit,backgrounds,calc,decorations.pathreplacing,patterns}

% Paleta institucional UEFS / Bioengenharia de Solos
\definecolor{bsNavy}{HTML}{1F4268}
\definecolor{bsNavyDark}{HTML}{152E4A}
\definecolor{bsGold}{HTML}{B38A3D}
\definecolor{bsGoldDark}{HTML}{8C6927}
\definecolor{bsTeal}{HTML}{0E7C7B}
\definecolor{bsForest}{HTML}{2D6A4F}
\definecolor{bsAlert}{HTML}{9E2A2B}
\definecolor{bsEarth}{HTML}{8C5A3C}
\definecolor{bsWater}{HTML}{2E86AB}
\definecolor{bsInk}{HTML}{26313D}
\definecolor{bsMuted}{HTML}{5B6573}

\definecolor{bgNavy}{HTML}{EDF2F7}
\definecolor{bgGold}{HTML}{F8F3E9}
\definecolor{bgTeal}{HTML}{EBF6F5}
\definecolor{bgForest}{HTML}{EEF5F0}
\definecolor{bgAlert}{HTML}{FDF2F2}

\tikzset{
  >=Stealth,
  bsBox/.style={
    draw=#1, fill=white, rounded corners=5pt, line width=1.15pt,
    inner sep=7pt, align=center, font=\footnotesize\sffamily, text=bsInk
  },
  bsTint/.style={
    draw=#1, fill=#1!8!white, rounded corners=5pt, line width=1.15pt,
    inner sep=7pt, align=center, font=\footnotesize\sffamily, text=bsInk
  },
  bsBadge/.style={
    fill=#1, text=white, rounded corners=9pt, inner xsep=8pt, inner ysep=3pt,
    font=\scriptsize\bfseries\sffamily
  },
  bsTitle/.style={font=\normalsize\bfseries\sffamily, text=bsNavy, align=center},
  bsSub/.style={font=\footnotesize\sffamily, text=bsMuted, align=center},
  bsArrow/.style={->, >=Stealth, line width=1.35pt},
  bsDArrow/.style={<->, >=Stealth, line width=1.35pt}
}

\tikzset{
  pill/.style={bsBadge=#1, rounded corners=4pt, minimum height=0pt, minimum width=0pt, text width=}
}

\begin{document}
\begin{tikzpicture}

% Painel esquerdo: Bacias de Captação / Barraginhas (bsWater)
\node[bsTint=bsWater, text width=6.15cm, minimum height=3.15cm, align=left, inner sep=10pt] (left) at (2.95, 2.35) {%
  \centering
  \tikz\node[pill=bsWater]{BACIAS DE CAPTAÇÃO / BARRAGINHAS};\\[8pt]
  \raggedright
  {\small\textbf{Posição na paisagem:} Terço médio e saídas laterais de estradas rurais}\\[5pt]
  {\small\textbf{Função hidráulica:} Armazenamento temporário de enxurrada e recarga do lençol freático}\\[5pt]
  {\small\textbf{Implantação:} Escavadas fora do talvegue profundo, com bordo dissipador de entrada}
};

% Painel direito: Paliçadas em Ravinas (bsEarth)
\node[bsTint=bsEarth, text width=6.15cm, minimum height=3.15cm, align=left, inner sep=10pt] (right) at (11.55, 2.35) {%
  \centering
  \tikz\node[pill=bsEarth]{PALIÇADAS EM RAVINAS};\\[8pt]
  \raggedright
  {\small\textbf{Posição na paisagem:} Dentro do canal erosivo ativo (ravina ou voçoroca)}\\[5pt]
  {\small\textbf{Função hidráulica:} Filtração do fluxo e formação de cunha de sedimentos}\\[5pt]
  {\small\textbf{Implantação:} Estruturas transversais em série para quebra de gradiente longitudinal}
};

% Seta bidirecional de integração com amplo espaço
\draw[bsDArrow, bsGoldDark] (left.east) -- (right.west);
\node[fill=bgGold, draw=bsGoldDark, rounded corners=4pt, inner xsep=6pt, inner ysep=3pt,
      font=\scriptsize\bfseries\sffamily, text=bsGoldDark, align=center]
  at (7.25, 2.35) {AÇÃO\\INTEGRADA};

% Divisor da faixa operacional com folga vertical clara
\draw[dashed, bsMuted!55, line width=0.9pt] (-0.50, -0.15) -- (15.00, -0.15);
\node[fill=white, inner xsep=10pt, font=\scriptsize\bfseries\sffamily, text=bsNavy]
  at (7.25, -0.15) {FLUXO OPERACIONAL EM ESTRADAS RURAIS};

% Faixa inferior: 4 estágios
\node[bsBox=bsNavy, text width=3.20cm, minimum height=2.25cm] (f1) at (1.25, -1.85) {%
  \tikz\node[pill=bsNavy]{1};\\[5pt]
  \textbf{Diagnóstico de Estrada\\\& Ponto de Saída}\\[3pt]
  {\scriptsize Mapeamento de bigodes\\e saídas de drenagem}
};

\node[bsBox=bsWater, text width=3.20cm, minimum height=2.25cm] (f2) at (5.25, -1.85) {%
  \tikz\node[pill=bsWater]{2};\\[5pt]
  \textbf{Escavação e\\Bordo Dissipador}\\[3pt]
  {\scriptsize Bacia semicircular e\\sangradouro protegido}
};

\node[bsBox=bsTeal, text width=3.20cm, minimum height=2.25cm] (f3) at (9.25, -1.85) {%
  \tikz\node[pill=bsTeal]{3};\\[5pt]
  \textbf{Decantação \&\\Infiltração}\\[3pt]
  {\scriptsize Retenção de finos e\\redução do pico de cheia}
};

\node[bsBox=bsForest, text width=3.20cm, minimum height=2.25cm] (f4) at (13.25, -1.85) {%
  \tikz\node[pill=bsForest]{4};\\[5pt]
  \textbf{Proteção da Borda e\\Manutenção Pós-Chuva}\\[3pt]
  {\scriptsize Gramíneas na crista e\\desassoreamento periódico}
};

\draw[bsArrow, bsNavy] (f1) -- (f2);
\draw[bsArrow, bsWater] (f2) -- (f3);
\draw[bsArrow, bsTeal] (f3) -- (f4);

% Título e subtítulo com folga acima dos painéis
\node[bsSub, anchor=south] (sub) at (7.25, 4.65)
  {Atuação complementar fora do talvegue (captação e recarga) e dentro do canal erosivo (filtração e retenção)};
\node[bsTitle, above=0.14cm of sub] (title)
  {Barraginhas (Bacias de Captação) vs. Paliçadas: Controle Integrado de Estradas Rurais};

\end{tikzpicture}
\end{document}
